% Generated draft: requires mathematical and source review.
\documentclass[11pt]{article}
\usepackage[margin=25mm]{geometry}
\usepackage{fontspec}
\setmainfont{Latin Modern Roman}
\usepackage{amsmath,amssymb,mathtools,unicode-math}
\setmathfont{Latin Modern Math}
\usepackage{enumitem,needspace,xurl,hyperref}
\setlength{\parindent}{0pt}
\setlength{\parskip}{5pt}
\setlength{\emergencystretch}{3em}
\newcommand{\N}{\mathbb{N}}
\newcommand{\Z}{\mathbb{Z}}
\newcommand{\Q}{\mathbb{Q}}
\newcommand{\R}{\mathbb{R}}
\newcommand{\C}{\mathbb{C}}
\newcommand{\F}{\mathbb{F}}
\newcommand{\A}{\mathbb{A}}
\newcommand{\PP}{\mathbb{P}}
\newcommand{\E}{\mathbb{E}}
\newcommand{\eps}{\varepsilon}
\newcommand{\dd}{\,\mathrm{d}}
\begin{document}
\section*{UnsolvedMath ID: 30006586}
\textbf{Source problem code:} catalogue entry
\subsection{Definitions and mathematical statement}
The supplied record concerns the report's algebras \(\mathcal G^{\mathrm f}\), \(\mathcal G\), and \(\mathcal Z_q\), and its primitive space \(\mathfrak{bm}_0\). Their defining generators, admissible indices, relations, coproduct, and grading are not supplied. Accordingly, the following is a statement of the intended conjectural targets in the report's notation, requiring those missing definitions before it becomes a self-contained mathematical problem.

\begin{enumerate}
\item Let \(\mathcal G^{\mathrm f}\) denote the intended formal algebra of balanced multiple \(q\)-zeta values modulo the balanced quasi-shuffle relations and \(\tau\)-invariance, and let \(\mathcal G\) denote the corresponding algebra of actual values. The conjecture is that these relations generate every algebraic relation among those values: the intended evaluation homomorphism
\[
\operatorname{ev}\colon\mathcal G^{\mathrm f}\longrightarrow\mathcal G
\]
is an algebra isomorphism, and
\[
\mathcal G^{\mathrm f}\cong\mathcal G=\mathcal Z_q.
\]
The precise admissible indices, meaning of actual \(q\)-zeta values, coefficient algebra, defining relations, and evaluation map must be supplied from the report. In particular, the excerpt does not specify whether \(q\) is formal or analytic, or an analytic domain if one is intended.

\item Equip the intended primitive space \(\mathfrak{bm}_0\) with the proposed bracket referred to as reference [5] in the excerpt. Prove that this bracket closes on \(\mathfrak{bm}_0\) and defines a Lie algebra. Thus, over its intended coefficient field, the bracket must be bilinear and satisfy
\[
[x,y]=-[y,x],\qquad
[x,[y,z]]+[y,[z,x]]+[z,[x,y]]=0
\quad(x,y,z\in\mathfrak{bm}_0).
\]
Neither the primitive condition nor the proposed bracket is defined in the supplied record, and reference [5] is not identified there.

\item Establish the proposed algebra isomorphism
\[
\mathcal G^{\mathrm f}\cong
\widetilde M_{\mathbb Q}\bigl(\operatorname{SL}_2(\mathbb Z)\bigr)
\otimes U(\mathfrak{bm}_0)^\vee.
\]
Here \(U(\mathfrak{bm}_0)\) denotes the universal enveloping algebra once the asserted Lie structure is established. The factor \(\widetilde M_{\mathbb Q}(\operatorname{SL}_2(\mathbb Z))\) conventionally denotes the algebra of quasimodular forms with rational coefficients, but its precise convention is not defined in the excerpt. The meaning of \(\vee\), the multiplication on the dual, the tensor-product conventions, and any required grading or completion are also unspecified. This factorization occurs in the supplied reconstructed statement but is absent from the supplied original excerpt.
\end{enumerate}

The original excerpt supplies additional qualifications and related objectives. It describes a term \(D(x)\) in an unspecified dimension conjecture as corresponding to a derivation of weight \(2\) on \(\mathfrak{bm}_0\). It attributes to Bachmann and van Ittersum, in its unidentified reference [3], the result that \(\mathcal G^{\mathrm f}\) is an \(\mathfrak{sl}_2\)-algebra, and in particular has a derivation \(\delta\) of weight \(-2\). For combinatorial multiple Eisenstein series with all entries \(k_i\geq2\), its displayed rule can only be transcribed with an unresolved coefficient:
\[
\delta\bigl(G(k_1,\ldots,k_d)\bigr)=
\begin{cases}
\kappa\,G(k_2,\ldots,k_d),&k_1=2,\\
0,&k_1\ne2.
\end{cases}
\]
The supplied extraction renders the coefficient as \texttt{-21}; its intended value \(\kappa\) cannot be recovered reliably from this record. The convention for an empty remaining index is also absent. The excerpt expresses an aim to relate \(\delta\) to ongoing joint work with Bachmann and van Ittersum, without giving a precise additional assertion.

Finally, it introduces the free associative algebra \(\mathbb Q\langle a,b,c\rangle\) and the anti-automorphism \(\tau\) specified by
\[
\tau(a)=c,\qquad\tau(b)=b,\qquad\tau(c)=a,\qquad
\tau(uv)=\tau(v)\tau(u).
\]
It then mentions an extension \(\Delta_c\) of a coproduct \(\Delta_b\), but the excerpt ends before defining it. No identification between this explicitly described \(\tau\) and the \(\tau\)-invariance in the first conjecture can be established from the supplied excerpt alone.
\subsection{Short English statement}
The intended conjecture says that balanced quasi-shuffle relations and \(\tau\)-invariance account for every relation among balanced multiple \(q\)-zeta values. It also proposes a Lie structure on \(\mathfrak{bm}_0\) and a factorization of the formal algebra using quasimodular forms and a dual universal enveloping algebra. Essential definitions and conventions are missing, so this is a qualified statement-only draft requiring human mathematical review.
\subsection{Sources}
\begin{itemize}
\item Local registry: \texttt{data/unsolvedmath/problems.json}, permanent UnsolvedMath ID \(30006586\), source problem code \texttt{catalogue entry}, title ``Lie Structure of Formal \(q\)-Zeta Algebras.'' The supplied record includes both a reconstructed statement and an incomplete original excerpt.
\item Supplied primary report citation: \emph{Multiple Zeta Values and Modular Forms} (2026), Oberwolfach Reports, DOI \texttt{10.4171/owr/2026/7}, pp.\ 433--486. \url{https://doi.org/10.4171/owr/2026/7}. This is also listed in the literature sources as ``Oberwolfach Report 7/2026.'' No individual proposer is identified in the record.
\item Supplied literature source: Bachmann--Kanno, \emph{Relations and Derivatives of Multiple Eisenstein Series}. \url{https://arxiv.org/abs/2602.08176}.
\item Literature assessment dated exactly \texttt{2026-08-22}: the supplied curation labels the problem open and states that the formal \(q\)-zeta algebra identification and the Lie-algebra assertion for \(\mathfrak{bm}_0\) remain conjectural in the 2026 report and related primary work. This is supplied dated curation, not an independently verified assessment; the record explicitly requests independent verification of time-sensitive status and cited claims.
\item The registry labels the statement \texttt{corrected\_verified} and reports that it was reconstructed and checked against the official Oberwolfach report. This is a supplied verification claim, not verification performed for this draft. The excerpt attributes the \(\mathfrak{sl}_2\)-algebra result to Bachmann and van Ittersum via [3] and the bracket formula to [5]; bibliographic identities for these numbered references are not supplied.
\end{itemize}
\end{document}
